\begin{thebibliography}{99}

\bibitem{ADM2019} V.~Alba, J.~Dubail and M.~Medenjak, Operator entanglement in interacting integrable quantum systems: the case of the rule 54 chain, \emph{Phys. Rev. Lett.} \textbf{122} (2019), 250603.

\bibitem{CaDiCaL} A.~Biere, K.~Fazekas, M.~Fleury and M.~Heisinger, CaDiCaL, Kissat, Paracooba, Plingeling and Treengeling entering the SAT Competition 2020, in \emph{Proc. SAT Competition 2020}, Univ. Helsinki, 2020, 51--53.

\bibitem{BT2000} F.~Blanchard and P.~Tisseur, Some properties of cellular automata with equicontinuity points, \emph{Ann. Inst. H. Poincar\'e Probab. Statist.} \textbf{36} (2000), 569--582.

\bibitem{BBGP1993} A.~Bobenko, M.~Bordemann, C.~Gunn and U.~Pinkall, On two integrable cellular automata, \emph{Comm. Math. Phys.} \textbf{158} (1993), 127--134.

\bibitem{BF1998} N.~Boccara and H.~Fuk\'s, Cellular automaton rules conserving the number of active sites, \emph{J. Phys. A} \textbf{31} (1998), 6007--6018.

\bibitem{BLR1988} M.~Boyle, D.~Lind and D.~Rudolph, The automorphism group of a shift of finite type, \emph{Trans. Amer. Math. Soc.} \textbf{306} (1988), 71--114.

\bibitem{BKP2021} B.~Bu\v{c}a, K.~Klobas and T.~Prosen, Rule 54: exactly solvable model of nonequilibrium statistical mechanics, \emph{J. Stat. Mech.} (2021), 074001.

\bibitem{CHR1979} E.~M.~Coven, G.~A.~Hedlund and F.~Rhodes, The commuting block maps problem, \emph{Trans. Amer. Math. Soc.} \textbf{249} (1979), 113--138.

\bibitem{CFK2019} V.~Cyr, J.~Franks and B.~Kra, The spacetime of a shift endomorphism, \emph{Trans. Amer. Math. Soc.} \textbf{371} (2019), 461--488.

\bibitem{CK2015} V.~Cyr and B.~Kra, The automorphism group of a shift of linear growth: beyond transitivity, \emph{Forum Math. Sigma} \textbf{3} (2015), e5.

\bibitem{PaperA} M.~Dayah, Reversible rule 110 has no hidden symmetries: six families of second-order cellular automata settled by edge moves, in preparation.

\bibitem{PaperB} M.~Dayah, Fronts at the speed of light: rigidity of permutive second-order cellular automata, in preparation.

\bibitem{PaperC} M.~Dayah, Every finite group from a wall: large automorphism groups of second-order cellular automata, in preparation.

\bibitem{PaperD} M.~Dayah, The 88 second-order elementary cellular automata: a classification by automorphism group, in preparation.

\bibitem{DDMP2016} S.~Donoso, F.~Durand, A.~Maass and S.~Petite, On automorphism groups of low complexity subshifts, \emph{Ergodic Theory Dynam. Systems} \textbf{36} (2016), 64--95.

\bibitem{MiniSat} N.~E\'en and N.~S\"orensson, An extensible SAT-solver, in \emph{Theory and Applications of Satisfiability Testing (SAT 2003)}, Lecture Notes in Comput. Sci. 2919, Springer, 2004, 502--518.

\bibitem{FNRW2021} P.~A.~Ferrari, C.~Nguyen, L.~T.~Rolla and M.~Wang, Soliton decomposition of the box-ball system, \emph{Forum Math. Sigma} \textbf{9} (2021), e60.

\bibitem{FGV2019} A.~J.~Friedman, S.~Gopalakrishnan and R.~Vasseur, Integrable many-body quantum Floquet--Thouless pumps, \emph{Phys. Rev. Lett.} \textbf{123} (2019), 170603.

\bibitem{Gopalakrishnan2018} S.~Gopalakrishnan, Operator growth and eigenstate entanglement in an interacting integrable Floquet system, \emph{Phys. Rev. B} \textbf{98} (2018), 060302(R).

\bibitem{GP2021} T.~Gombor and B.~Pozsgay, Integrable spin chains and cellular automata with medium-range interaction, \emph{Phys. Rev. E} \textbf{104} (2021), 054123.

\bibitem{GP2022} T.~Gombor and B.~Pozsgay, Superintegrable cellular automata and dual unitary gates from Yang--Baxter maps, \emph{SciPost Phys.} \textbf{12} (2022), 102.

\bibitem{GP2024} T.~Gombor and B.~Pozsgay, Integrable deformations of superintegrable quantum circuits, \emph{SciPost Phys.} \textbf{16} (2024), 114.

\bibitem{GHKV2018} S.~Gopalakrishnan, D.~A.~Huse, V.~Khemani and R.~Vasseur, Hydrodynamics of operator spreading and quasiparticle diffusion in interacting integrable systems, \emph{Phys. Rev. B} \textbf{98} (2018), 220303(R).

\bibitem{HT1991} T.~Hattori and S.~Takesue, Additive conserved quantities in discrete-time lattice dynamical systems, \emph{Phys. D} \textbf{49} (1991), 295--322.

\bibitem{Hedlund1969} G.~A.~Hedlund, Endomorphisms and automorphisms of the shift dynamical system, \emph{Math. Systems Theory} \textbf{3} (1969), 320--375.

\bibitem{Hochman2010} M.~Hochman, On the automorphism groups of multidimensional shifts of finite type, \emph{Ergodic Theory Dynam. Systems} \textbf{30} (2010), 809--840.

\bibitem{IT2018} A.~Inoue and S.~Takesue, Two extensions of exact nonequilibrium steady states of a boundary-driven cellular automaton, \emph{J. Phys. A} \textbf{51} (2018), 425001.

\bibitem{PySAT} A.~Ignatiev, A.~Morgado and J.~Marques-Silva, PySAT: a Python toolkit for prototyping with SAT oracles, in \emph{Theory and Applications of Satisfiability Testing (SAT 2018)}, Lecture Notes in Comput. Sci. 10929, Springer, 2018, 428--437.

\bibitem{KT2015} J.~Kari and S.~Taati, Statistical mechanics of surjective cellular automata, \emph{J. Stat. Phys.} \textbf{160} (2015), 1198--1243.

\bibitem{KS2000} B.~Kitchens and K.~Schmidt, Isomorphism rigidity of irreducible algebraic $\mathbb{Z}^d$-actions, \emph{Invent. Math.} \textbf{142} (2000), 559--577.

\bibitem{KlobasThesis} K.~Klobas, \emph{Exact time-dependent solutions of interacting systems}, Doctoral thesis, University of Ljubljana, 2020.

\bibitem{Kopra2020} J.~Kopra, Glider automorphisms and a finitary Ryan's theorem for transitive subshifts of finite type, \emph{Nat. Comput.} \textbf{19} (2020), 773--786.

\bibitem{Kopra2022} J.~Kopra, Glider automata on all transitive sofic shifts, \emph{Ergodic Theory Dynam. Systems} \textbf{42} (2022), 3716--3744.

\bibitem{Kurka1997} P.~K\r{u}rka, Languages, equicontinuity and attractors in cellular automata, \emph{Ergodic Theory Dynam. Systems} \textbf{17} (1997), 417--433.

\bibitem{LM1995} D.~Lind and B.~Marcus, \emph{An Introduction to Symbolic Dynamics and Coding}, Cambridge University Press, Cambridge, 1995.

\bibitem{MBG2004} A.~Moreira, N.~Boccara and E.~Goles, On conservative and monotone one-dimensional cellular automata and their particle representation, \emph{Theoret. Comput. Sci.} \textbf{325} (2004), 285--316.

\bibitem{PP2026} C.~Paletta and T.~Prosen, On the integrability structure of the deformed rule-54 reversible cellular automaton, preprint, arXiv:2603.25424, 2026.

\bibitem{Pivato2002} M.~Pivato, Conservation laws in cellular automata, \emph{Nonlinearity} \textbf{15} (2002), 1781--1793.

\bibitem{Ryan1972} J.~P.~Ryan, The shift and commutativity, \emph{Math. Systems Theory} \textbf{6} (1972), 82--85.

\bibitem{Salo2019} V.~Salo, Transitive action on finite points of a full shift and a finitary Ryan's theorem, \emph{Ergodic Theory Dynam. Systems} \textbf{39} (2019), 1637--1667.

\bibitem{Takesue1987} S.~Takesue, Reversible cellular automata and statistical mechanics, \emph{Phys. Rev. Lett.} \textbf{59} (1987), 2499--2502.

\bibitem{Takesue1989} S.~Takesue, Ergodic properties and thermodynamic behavior of elementary reversible cellular automata. I. Basic properties, \emph{J. Stat. Phys.} \textbf{56} (1989), 371--402.

\bibitem{Takesue1995} S.~Takesue, Staggered invariants in cellular automata, \emph{Complex Systems} \textbf{9} (1995), 149--168.

\bibitem{Takesue1997} S.~Takesue, Seru \=otomaton no hozonry\=o [Conserved quantities of cellular automata] (in Japanese), \emph{RIMS K\^oky\^uroku} \textbf{1020} (1997), 103--126.

\bibitem{TM1990} T.~Toffoli and N.~Margolus, Invertible cellular automata: a review, \emph{Phys. D} \textbf{45} (1990), 229--253.

\bibitem{Wolfram2002} S.~Wolfram, \emph{A New Kind of Science}, Wolfram Media, Champaign, IL, 2002.

\end{thebibliography}
